\section{Related Work}
\label{sec:related_work}

\para{Number representations in LLMs.}
A growing body of work investigates how large language models encode and process numbers. \citet{levy2024language} demonstrated that LLMs represent numbers using per-digit circular encodings in base 10, achieving over 90\% accuracy with linear probes on models including \mistral and Llama~3. They showed that causal interventions on these digit representations modify model outputs, confirming that the representations are functionally relevant. \citet{kadlcik2025pretrained} extended these findings across multiple model families (Llama~3, Phi~4, OLMo~2) and scales (1B--72B parameters), showing that sinusoidal probes achieve near-perfect accuracy and that embedding precision explains a large fraction of arithmetic errors. \citet{stefanik2025unravelling} showed that different LLMs learn interchangeable, universal number representations that generalize across hidden states and input contexts. We build directly on the circular probing methodology of \citet{levy2024language}, extending it from digit strings and preliminary English word-form results to a systematic cross-lingual comparison involving French vigesimal number words.

\para{Tokenization and numerical reasoning.}
How numbers are tokenized strongly influences model performance on arithmetic tasks. \citet{singh2024tokenization} showed that right-to-left tokenization improves arithmetic accuracy and that model errors follow systematic patterns suggesting structured computation. \citet{nogueira2021investigating} found that subword tokenization fails for multi-digit addition, while position-augmented formats can enable accurate computation. \citet{zhou2024scaling} studied how different numeral systems (base 10, 100, 1000) affect learning efficiency, finding that base 10 is consistently more data-efficient. French vigesimal number words present an extreme case of ``unusual surface form,'' where the tokenized representation encodes arithmetic operations rather than digit values. Our work complements these studies by examining how tokenization of complex number words affects the quality of internal numeric representations rather than output accuracy.

\para{Numerical reasoning and arithmetic in LLMs.}
The broader landscape of numerical reasoning in LLMs includes work on arithmetic circuits \citep{hanna2023does}, the dissociation between language and arithmetic representations \citep{kisako2025representational}, and the impact of semantic cues on computation \citep{deleeuw2025semantic}. \citet{hanna2023does} identified specific circuits in GPT-2 that compute greater-than comparisons, revealing that MLPs boost probabilities of numerically valid continuations. \citet{kisako2025representational} showed that arithmetic equations and general language occupy separate regions in LLM representation space, raising the question of whether French number words activate the ``language'' region, the ``arithmetic'' region, or both. \citet{marjieh2025number} used similarity-based prompting to show that LLM number representations blend string-like (Levenshtein distance) and numerical (log-linear distance) structure. Our work provides a new lens on these questions: French vigesimal numbers are simultaneously linguistic (word-form) and arithmetic (encoding $4\times20+10+8$).

\para{Cross-lingual and multilingual representations.}
\citet{zada2025brains} showed that language models trained on different languages converge onto shared conceptual spaces, particularly in middle layers, and that encoding models trained on one language can predict neural activity in listeners of other languages. \citet{xu2024numcot} evaluated LLMs on numeral conversions across different counting systems, finding that non-standard numeral systems present significant challenges. Our work leverages the French--Belgian French contrast as a within-language controlled comparison: the same numbers, in the same language family, but with vigesimal versus decimal structure. This allows us to isolate the effect of counting system architecture on internal representations while controlling for language-level confounds.
